\scanpage{94}\mspage{85}
Quand le magma est-il fidèle \unsure{?}\,, I.e. quand a-t-on la condition de fidélité
\begingroup
\makeatletter\tagsleft@true\makeatother
\[
 \mathrm{Omb}^{\circ}(X)\cap\mathrm{Omb}^{\circ}(Y)=\varnothing
 \Longrightarrow X\,|\circ|\,Y\quad ? \tag*{(*)}
\]
\endgroup

Il faut déjà que sur $L$
\[
 P_x\ne P_y\Longrightarrow x\,|\circ|\,y
\]
i.e. que $L$ soit \uline{totalement} ordonné. Cela équivaut à $(*)$ pour $X=P_x$, $Y=P_y$. Dans le cas $X=P_x$, $Y=S_{a,b}$, cela signifie $x\notin ]a,b[\Longrightarrow x\,|\circ|\,S_{ab}$ i.e. $x\leq a$ ou $x\geq b$, OK. Dans le cas où $\deleted{\unsure{x}}X=S_{a,b}$, $Y=S_{a',b'}$, cela signifie que l'on aura \deleted{$S_{a,b}$} si on \deleted{$S_{a,b}$} n'a pas $b\leq\unsure{a'|b'}$ ou $b'\leq a$, i.e. si on a à la fois $\underline{b>a'}$, $\underline{b'>a}$, alors il y a un \unsure{raff} interne commun. Je distingue quatre cas, suivant la position relative de $a$ p.r. à $a'$, et de $b$ p.r. à $b'$

\[
\begin{aligned}
 {}\cdot\quad a'\leq a<b\leq b'&\quad\text{on prend }\deleted{\unsure{[}}S_{ab}\ \text{comme}\\
 &\quad\text{raffinement interne commun}\\[3pt]
 a\leq a'\unsure{<|\leq}b\leq b'\quad\text{ou}\quad a'\leq a<b'\leq b
 &\quad\text{on prend }S_{a'b}\\
 &\qquad\text{resp }S_{ab'}\\[3pt]
 a\leq a'\unsure{<|\leq}b'\leq b&\quad\text{on prend }S_{a',b'}
\end{aligned}
\]

\unsure{Donc|D'où} \uline{Lemme}\qquad $M$ est fidèle sss $L$ est totalement ordonné

\uline{Corollaire}\qquad $M$ est ponctuel sss $L$ est totalement ordonné et divisible (et alors $M$ est \unsure{$=$|un} ens \unsure{multistr.}).
% archive sheet number: 94

